\documentclass[12pt,a4paper]{article}
\usepackage[margin=25mm, headheight=15pt, headsep=10pt]{geometry}
\usepackage{fontspec}
\usepackage[table]{xcolor} % Note: table option is required for rowcolors
\usepackage{titlesec}
\usepackage{tocloft}
\usepackage{fancyhdr}
\usepackage{booktabs}
\usepackage{tcolorbox}
\tcbuselibrary{skins, breakable}
\usepackage[colorlinks=true, linkcolor=emerald, urlcolor=emerald, bookmarks=true]{hyperref}

\definecolor{emerald}{RGB}{0,102,51}
\definecolor{darkred}{RGB}{139,0,0}
\definecolor{lightgray}{RGB}{245,245,245}

\usepackage[nil]{babel}
\babelprovide[import, main]{bengali}
\babelprovide[import]{arabic}
\babelfont{rm}[Renderer=HarfBuzz,Scale=1.0]{Noto Serif Bengali}
\babelfont{sf}[Renderer=HarfBuzz]{Noto Sans Bengali}
\babelfont{tt}[Renderer=HarfBuzz]{Noto Sans Bengali}
\babelfont[arabic]{rm}[Renderer=HarfBuzz,Scale=1.1]{Noto Naskh Arabic}

% Section styling
\titleformat{\section}{\Large\bfseries\color{emerald}}{\thesection.}{0.5em}{}
\titleformat{\subsection}{\large\bfseries\color{emerald!85!black}}{\thesubsection.}{0.5em}{}
\titleformat{\subsubsection}{\normalsize\bfseries\color{emerald!70!black}}{\thesubsubsection.}{0.5em}{}
\titlespacing*{\section}{0pt}{18pt}{8pt}

% Fancy Header/Footer setup
\pagestyle{fancy}
\fancyhf{} % clear all header and footer fields
\fancyhead[L]{\color{emerald}\small\textbf{\nouppercase{\leftmark}}}
\fancyfoot[L]{\scriptsize\color{black!50}{\lat{AI}}-সংকলিত, আলেম-যাচাইকৃত নয়}
\fancyfoot[C]{\color{emerald!70!black}\thepage}
\renewcommand{\headrulewidth}{0.4pt}
\renewcommand{\headrule}{\hbox to\headwidth{\color{emerald}\leaders\hrule height \headrulewidth\hfill}}
\renewcommand{\footrulewidth}{0pt}

% Custom boxes
% 1. Ayat/Hadith Box
\newtcolorbox{ayatbox}[1][]{
    enhanced, breakable, 
    colback=emerald!5, 
    colframe=emerald, 
    boxrule=0.5pt, 
    arc=3pt, 
    left=10pt, right=10pt, top=10pt, bottom=10pt,
    #1
}

% 2. Quote/Translation Box
\newtcolorbox{quotebox}[1][]{
    enhanced, breakable,
    frame hidden,
    colback=lightgray,
    borderline west={3pt}{0pt}{emerald},
    left=15pt, right=10pt, top=10pt, bottom=10pt,
    sharp corners,
    #1
}

% 3. AI-generated notice box
\newtcolorbox{warnbox}[1][]{
    enhanced, breakable,
    colback=red!4,
    colframe=darkred,
    boxrule=0.8pt,
    arc=3pt,
    left=10pt, right=10pt, top=10pt, bottom=10pt,
    #1
}

% Commands
\newcommand{\arab}[1]{\foreignlanguage{arabic}{#1}}
\newcommand{\kib}[1]{\textcolor{emerald}{\textbf{#1}}}

% Latin (English) fallback: Noto Serif Bengali has NO Latin glyphs
\newfontfamily\latfont[Renderer=HarfBuzz]{Noto Serif}
\newcommand{\lat}[1]{{\latfont #1}}

\setlength{\parskip}{5pt}
\setlength{\parindent}{0pt}

% Fix for missing bullet in Noto Serif Bengali
\renewcommand{\labelitemi}{$\bullet$}



\begin{document}
\begin{center}{\Large\bfseries\color{emerald} আল্লাহর রহমত ও তাওবা — \lat{03}-মায়েয-ও-গামেদিয়া-পাপ-স্বীকার-ও-ক্ষমা}\end{center}
\vspace{1em}
\section*{ঘটনা ৩ — মায়েয ও গামেদিয়া: পাপ স্বীকার, শাস্তি মেনে নেওয়া ও নবীজির ঘোষণা}

\begin{warnbox}
 \textbf{গুরুত্বপূর্ণ নোটিশ (\lat{Important} \lat{Notice}):} এই কন্টেন্টটি \lat{AI} (কৃত্রিম বুদ্ধিমত্তা) দিয়ে প্রস্তুত ও সংকলিত। \lat{AI} কঠোর সূত্র-মান নীতি মেনে শুধু নির্ভরযোগ্য উৎস থেকে তথ্য সংগ্রহ করেছে — কেবল সহীহ/হাসান হাদিস ও সহীহ সনদের আসার রাখা হয়েছে, দুর্বল সনদ বাদ দেওয়া হয়েছে। তবে এটি কোনো আলেম বা মুহাদ্দিস কর্তৃক যাচাইকৃত নয়।
\end{warnbox}

\begin{quote}
\textbf{সম্পর্কিত ফাইল:} হাদিস ফাইল — যিনার শাস্তি ও তাওবা; ব্যবহারিক গাইড — মানুষের হক ও তাওবার শর্ত।
\end{quote}

\medskip\hrule\medskip

\subsection*{১. সূত্র কার্ড (\lat{Source} \lat{Card})}

\begin{table}[h]\centering\small
\begin{tabular}{ll}
বিষয় & বিবরণ \\
\hline
\textbf{ঘটনা} & মায়েয ইবনে মালিক (রাঃ) ও গামেদিয়া (জুহাইনা) নারীর পাপ স্বীকার — হদ (রজম) আদায় — নবীজির (সা.) কাছে তাদের তাওবার সর্বোচ্চ প্রশংসা \\
\textbf{বর্ণনাকারী} & বুরাইদা ইবনে হুসাইব (রাঃ); ইমরান ইবনে হুসাইন (রাঃ) \\
\textbf{গ্রন্থ ও নম্বর} & সহীহ মুসলিম ১৬৯৫, ১৬৯৬; সহীহ বুখারি ৬৮২৭, ৬৮২৮ \\
\textbf{অধ্যায়} & কিতাবুল হুদুদ — বাবু মানি'তারাফা আলা নাফসিহি বিজ যিনা (যে নিজের যিনা স্বীকার করে) \\
\textbf{সনদের মান} & \textbf{সহীহ} (বুখারি-মুসলিম) \\
\textbf{স্থান ও সময়} & মদিনা — নবীজির (সা.) জীবদ্দশায় \\
\hline
\end{tabular}
\end{table}

\medskip\hrule\medskip

\subsection*{২. সংক্ষিপ্ত পটভূমি (\lat{Background})}

ইসলামে পাপ গোপন রাখা ও আল্লাহর কাছে তাওবা করা সাধারণ নিয়ম। কিন্তু মায়েয ও গামেদিয়া নিজেরা এসে পাপ স্বীকার করেছিলেন এবং শাস্তি (হদ) মেনে নিয়েছিলেন — শাস্তির মাধ্যমে নিজেদের পবিত্র করতে চেয়েছিলেন। নবীজি (সা.) চারবার মায়েযকে ফিরিয়ে দিয়েছিলেন — যেন সে তাওবা করে চলে যায়; কিন্তু সে জিদ ধরে রইল। এই ঘটনা দুটিই দেখায় — পাপ স্বীকার ও শাস্তি মেনে নেওয়ার অন্তর আল্লাহর কাছে কত মূল্যবান।

\medskip\hrule\medskip

\subsection*{৩. পূর্ণ ঘটনার বিশুদ্ধ বর্ণনা (\lat{The} \lat{Full} \lat{Narration})}

\subsubsection*{পর্ব ১ — মায়েয ইবনে মালিক}

মায়েয ইবনে মালিক নবীজির (সা.) কাছে এসে বললেন: \textbf{\lat{“}হে আল্লাহর রাসূল! আমাকে পবিত্র করুন।\lat{”}} নবীজি (সা.) বললেন: \textbf{\lat{“}তোমার দুর্ভোগ! ফিরে যাও — আল্লাহর কাছে ক্ষমা চাও এবং তাওবা করো।\lat{”}} সে ফিরে গেল, কিন্তু বেশি দূরে যায়নি — আবার এসে বলল: আমাকে পবিত্র করুন। এভাবে চারবার। চতুর্থবার নবীজি (সা.) জিজ্ঞেস করলেন: \textbf{\lat{“}কিসের থেকে পবিত্র করব?\lat{”}} সে বলল — যিনা থেকে। নবীজি (সা.) জিজ্ঞেস করলেন — সে কি পাগল? জানানো হলো — না। জিজ্ঞেস করলেন — সে কি মদ খেয়েছে? এক ব্যক্তি দাঁড়িয়ে তার নিঃশ্বাস শুঁকলেন — মদের গন্ধ নেই। নবীজি (সা.) জিজ্ঞেস করলেন: \textbf{\lat{“}তুমি কি যিনা করেছ?\lat{”}} সে বলল — হ্যাঁ। নবীজি (সা.) তার ব্যাপারে আদেশ দিলেন — তাকে পাথর মারা হলো।

লোকজন দুই দলে ভাগ হলো — এক দল বলল: সে ধ্বংস হয়ে গেছে, তার পাপ তাকে ঘিরে ফেলেছে; আরেক দল বলল: মায়েযের চেয়ে উত্তম তাওবা আর নেই — সে নবীজির (সা.) কাছে এসে তার হাতে হাত রেখে বলেছিল: পাথর দিয়ে আমাকে হত্যা করো। দুই-তিন দিন এই বিতর্ক চলল। এরপর নবীজি (সা.) বসে থাকা সাহাবিদের কাছে এসে সালাম দিয়ে বসলেন এবং বললেন: \textbf{\lat{“}মায়েয ইবনে মালিকের জন্য ক্ষমা প্রার্থনা করো।\lat{”}} সাহাবিরা বললেন — আল্লাহ মায়েযকে ক্ষমা করুন। নবীজি (সা.) বললেন: \textbf{\lat{“}সে এমন তাওবা করেছে, যা একটি সম্প্রদায়ের মধ্যে বণ্টন করা হলে তাদের সকলের জন্য যথেষ্ট হতো।\lat{”}}

\subsubsection*{পর্ব ২ — গামেদিয়া (জুহাইনা) নারী}

এরপর গামেদ গোত্রের (অন্য বর্ণনায় জুহাইনা গোত্রের) এক নারী নবীজির (সা.) কাছে এসে বলল: আমাকে পবিত্র করুন। নবীজি (সা.) বললেন: \textbf{\lat{“}তোমার দুর্ভোগ! ফিরে যাও — আল্লাহর কাছে ক্ষমা চাও এবং তাওবা করো।\lat{”}} সে বলল: আমি দেখছি, আপনি আমাকে ফিরিয়ে দিতে চান, যেমন মায়েযকে ফিরিয়ে দিয়েছিলেন। নবীজি (সা.) জিজ্ঞেস করলেন: \textbf{\lat{“}তোমার কী হয়েছে?\lat{”}} সে বলল — যিনার কারণে আমি গর্ভবতী। নবীজি (সা.) বললেন: \textbf{\lat{“}তুমি?\lat{”}} সে বলল — হ্যাঁ। নবীজি (সা.) বললেন: \textbf{(শাস্তি হবে না) যতক্ষণ না তুমি গর্ভের সন্তান প্রসব করো।} এক আনসারি তার দায়িত্ব নিলেন — সন্তান প্রসবের পর তাকে নবীজির (সা.) কাছে আনা হলো। নবীজি (সা.) বললেন: \textbf{\lat{“}এখন আমরা তাকে পাথর মারব না এবং তার শিশুকে দুধপান করানোর কেউ থাকবে না।\lat{”}} এক আনসারি উঠে বললেন — শিশুর দুধপান করার দায়িত্ব আমি নিচ্ছি। এরপর নবীজি (সা.) তার ব্যাপারে আদেশ দিলেন — তার কাপড় বাঁধা হলো, তারপর তাকে পাথর মারা হলো। এরপর নবীজি (সা.) তার জানাজার নামাজ পড়ালেন। উমার (রাঃ) বললেন: হে আল্লাহর রাসূল! আপনি তার জানাজা পড়ালেন — অথচ সে যিনা করেছে! নবীজি (সা.) বললেন: \textbf{\lat{“}সে এমন তাওবা করেছে, যা মদিনার সত্তর জনের মধ্যে বণ্টন করা হলে তাদের জন্য যথেষ্ট হতো। তুমি কি এমন তাওবা দেখেছ — এর চেয়ে উত্তম, যে নিজের প্রাণ আল্লাহর জন্য উৎসর্গ করল?\lat{”}}

\medskip\hrule\medskip

\subsection*{৪. শব্দের অর্থ (\lat{Key} \lat{Words})}

\begin{table}[h]\centering\small
\begin{tabular}{ll}
শব্দ & অর্থ \\
\hline
\textbf{ত্বাহহিরনী (\arab{طهرني})} & আমাকে পবিত্র করুন — শাস্তির মাধ্যমে পরিশুদ্ধ করার আবেদন \\
\textbf{ওয়াইলাকা (\arab{ويلك})} & তোমার দুর্ভোগ — স্নেহ ও করুণার উক্তি \\
\textbf{রুজিমা (\arab{رجم})} & পাথর মারা — হদ (যিনার শাস্তি) \\
\textbf{লাকাদ তাবাত তাওবাতান} & সে এমন তাওবা করেছে \\
\textbf{জাদাত বিনাফসিহা লিল্লাহ (\arab{جادت} \arab{بنفسها} \arab{لله})} & নিজের প্রাণ আল্লাহর জন্য বিলিয়ে দিয়েছে \\
\hline
\end{tabular}
\end{table}

\medskip\hrule\medskip

\subsection*{৫. সংশ্লিষ্ট দলিল ও রেফারেন্স}

\begin{quote}
\textbf{\lat{“}সে এমন তাওবা করেছে, যা একটি সম্প্রদায়ের মধ্যে বণ্টন করা হলে সকলের জন্য যথেষ্ট হতো।\lat{”}} — সহীহ মুসলিম ১৬৯৫
\end{quote}

\begin{quote}
\textbf{\lat{“}সে এমন তাওবা করেছে, যা মদিনার সত্তর জনের মধ্যে বণ্টন করা হলে তাদের জন্য যথেষ্ট হতো… তুমি কি এমন তাওবা দেখেছ, যে নিজের প্রাণ আল্লাহর জন্য উৎসর্গ করল?\lat{”}} — সহীহ মুসলিম ১৬৯৬
\end{quote}

\begin{quote}
\textbf{\lat{“}বলো: হে আমার বান্দারা… তোমরা আল্লাহর রহমত থেকে নিরাশ হয়ো না।\lat{”}} — সূরা আয-যুমার ৩৯:৫৩
\end{quote}

\subsubsection*{মূল শিক্ষা (দলিলভিত্তিক)}

\begin{enumerate}
\item \textbf{পাপ স্বীকার ও হদ — তাওবার চূড়া:} নবীজি (সা.) স্বয়ং ঘোষণা করেছেন — এমন তাওবা সম্প্রদায়/সত্তর জনের জন্য যথেষ্ট; অন্তরের আন্তরিক ফেরা পাপের হিসাব মুছে দেয়।
\item \textbf{নবীজির (সা.) পদ্ধতি:} বারবার ফিরিয়ে দেওয়া — যেন তাওবা করে চলে যায়; শাস্তি দিতে নবীজির (সা.) তাড়া ছিল না — কিন্তু বান্দা যখন জিদ ধরে, তখন হদ আদায়ই আখিরাতের শাস্তি থেকে মুক্তির পথ।
\item \textbf{লোকের বিচার ভুল হতে পারে:} সাহাবিরাও দ্বিধাবিভক্ত ছিলেন — কেউ ভেবেছিলেন মায়েয ধ্বংস হয়েছে; আল্লাহ ও তাঁর রাসূলের সিদ্ধান্তই চূড়ান্ত। পাপীর ভবিষ্যৎ আল্লাহর হাতে — মানুষের ধারণায় নয়।
\item \textbf{গোপন রাখার নিয়ম প্রসঙ্গ:} সাধারণ নিয়ম — পাপ গোপন রাখা ও আল্লাহর কাছে তাওবা; তবে হক আদায় ও হদ (আদালতের ক্ষেত্রে) আলাদা বিধান — এ ঘটনায় স্বেচ্ছা স্বীকারোক্তির মর্যাদা দেখানো হয়েছে।
\end{enumerate}
\medskip\hrule\medskip

\subsection*{৬. রেফারেন্স}

\begin{itemize}
\item সহীহ মুসলিম ১৬৯৫ (মায়েয), ১৬৯৬ (গামেদিয়া/জুহাইনা)
\item সহীহ বুখারি ৬৮২৭, ৬৮২৮
\item কুরআন: যুমার ৩৯:৫৩, ফুরকান ২৫:৭০
\item ব্যাখ্যা: শারহু নববী — হদ ও তাওবার আলোচনা; ফাতহুল বারি
\end{itemize}

\end{document}
